\chapter{Logical structure}
\chaptercontents{A & Propositional logic (\S1.1) \\ B & Variables and quantifiers (\S1.2) \\ C & Logical equivalence (\S1.3) \\ D & Chapter 1 exercises (1.E)}%
{\fE\ 1.1: 21, 25, 32, 33, 38, 47, 56, 57 \textperiodcentered\ 1.2: 6, 15, 18, 39 \textperiodcentered\ 1.3: 5, 13, 14, 22, 27, 44ace \textperiodcentered\ 1.E: 18\\ \fR\ 1.1: 37, 43, 54, 55, 65, 66 \textperiodcentered\ 1.2: 21, 22, 28, 37ab \textperiodcentered\ 1.3: 32, 35 \textperiodcentered\ 1.E: 11, 28, 29, 30\\ Bonus: 1.2.14. \S1.1 is in \textbf{Quiz 1}.}

The chapter's one idea: a proposition's \emph{logical structure} (which operators and quantifiers it is built from) tells you what to write. Each operator has \textbf{introduction rules} (how to prove a goal of that shape) and \textbf{elimination rules} (how to use an assumption of that shape); each rule is a \textbf{Strategy}.

\hsection{Propositional logic}

\begin{key}[{Example 1.1.1 \textperiodcentered\ assumptions and goals}]
At every point of a proof you have a list of \term{assumptions} (what is known or supposed) and \term{goals} (what remains to be shown). Every sentence you write changes the lists: ``Suppose $P$'' adds $P$ to the assumptions; unfolding a definition (``$b=qc$ for some $q\in\Z$'') replaces an assumption by its definition and introduces a new variable; ``it suffices to show \dots'' replaces the goal. Keep both lists in your head (or on scrap paper) at all times.
\end{key}

\begin{key}[{Definitions 1.1.2--1.1.3}]
A \term{propositional variable} ($p,q,r,s$) is a symbol representing a proposition. A \term{propositional formula} is a propositional variable, or is built from simpler formulae using \term{logical operators} ($\wedge,\vee,\Rightarrow,\neg,\Leftrightarrow$). Its truth value is determined by the truth values of its subformulae.
\end{key}

\hsub{Conjunction $\wedge$ (``and'')}
\begin{key}[{Definition 1.1.4 \textperiodcentered\ rules for $\wedge$}]
($\wedge$I) if $p$ and $q$ are true then $p\wedge q$ is true;\quad ($\wedge$E$_1$) if $p\wedge q$ is true then $p$ is true;\quad ($\wedge$E$_2$) if $p\wedge q$ is true then $q$ is true. Watch for hidden ``and''s: ``$7$ is a prime factor of $28$'' is ``$7$ is prime'' $\wedge$ ``$7$ divides $28$''.
\end{key}
\begin{strategy}[{Strategies 1.1.7 and 1.1.10}]
\textbf{Proving $p\wedge q$:} prove $p$, then prove $q$ (two proofs tied together). \textbf{Assuming $p\wedge q$:} you may assume $p$ and assume $q$ separately.
\end{strategy}
Rules are written in natural-deduction style, premises above a line and conclusion below; stacking rules gives a \term{proof tree}, which is a plan for a proof (page 36).

\begin{trybox}[{1.1.6, 1.1.13, 1.1.14 \fO}]
\textbf{1.1.6} Express ``John is a mathematician who lives in Pittsburgh'' as $p\wedge q$.\quad \textbf{1.1.13} Given $(p\wedge q)\wedge r$, prove $q$ using Strategy 1.1.10.\quad \textbf{1.1.14} Proof tree for $(p\wedge q)\wedge(r\wedge s)$; apply it to ``$2$ is an even prime number and $3$ is an odd prime number''.
\end{trybox}

\hsub{Disjunction $\vee$ (``or'')}
\begin{key}[{Definition 1.1.15 \textperiodcentered\ rules for $\vee$}]
($\vee$I$_1$) if $p$ is true then $p\vee q$ is true;\quad ($\vee$I$_2$) if $q$ is true then $p\vee q$ is true;\quad ($\vee$E) if $p\vee q$ is true, and $r$ can be derived from $p$ and also from $q$, then $r$ is true. (``Or'' is inclusive.)
\end{key}
\begin{strategy}[{Strategies 1.1.16 and 1.1.19}]
\textbf{Proving $p\vee q$:} it suffices to prove just one of $p$ or $q$. \textbf{Assuming $p\vee q$ (proof by cases):} to derive $r$, first derive $r$ from the temporary assumption $p$, then derive $r$ from the temporary assumption $q$. Write ``Case 1 \dots\ Case 2 \dots\ In both cases, $r$.''
\end{strategy}
The division theorem produces disjunctions to split on: $n=3k$ or $n=3k+1$ or $n=3k+2$ (Proposition 1.1.24: $n^2$ leaves remainder $0$ or $1$ on division by $3$).

\begin{bookex}[essential]{1.1.21}
Let $p,q,r$ be propositions and suppose $q\vee r$ is true. Prove that $(p\vee q)\vee(p\vee r)$ is true, using Strategies 1.1.19 and 1.1.16.
\solution
We split into cases on the assumption $q\vee r$ (Strategy 1.1.19).\\
\emph{Case 1:} $q$ is true. Then $p\vee q$ is true by Strategy 1.1.16 (rule $\vee$I$_2$), and hence $(p\vee q)\vee(p\vee r)$ is true by Strategy 1.1.16 ($\vee$I$_1$).\\
\emph{Case 2:} $r$ is true. Then $p\vee r$ is true ($\vee$I$_2$), and hence $(p\vee q)\vee(p\vee r)$ is true ($\vee$I$_2$).\\
In both cases $(p\vee q)\vee(p\vee r)$ is true.\qed
\end{bookex}

\begin{bookex}[essential]{1.1.25}
Let $n$ be an integer. Prove that $n^2$ leaves a remainder of $0$, $1$ or $4$ when divided by $5$.
\solution
By the division theorem (Theorem 0.42), $n=5k+r$ for some $k\in\Z$ and $r\in\{0,1,2,3,4\}$. Then
\[
n^2=25k^2+10kr+r^2=5(5k^2+2kr)+r^2 .
\]
We split into cases on $r$. If $r=0$: $n^2=5(5k^2)$, remainder $0$. If $r=1$: $n^2=5(5k^2+2k)+1$, remainder $1$. If $r=2$: $n^2=5(5k^2+4k)+4$, remainder $4$. If $r=3$: $r^2=9=5+4$, so $n^2=5(5k^2+6k+1)+4$, remainder $4$. If $r=4$: $r^2=16=15+1$, so $n^2=5(5k^2+8k+3)+1$, remainder $1$.\\
In all cases the number in brackets is an integer and the remainder is $0$, $1$ or $4$ (uniqueness in Theorem 0.42).\qed
\end{bookex}

\begin{trybox}[{1.1.23, 1.1.26 \fO}]
\textbf{1.1.23} Let $n$ be $8$ or $9$. Prove that $n$ is even or a perfect square.\quad \textbf{1.1.26} Let $a,b\in\R$ with $a^2-4b\ne0$ and let $\alpha,\beta$ be the roots of $x^2+ax+b$. Prove there is $c\in\R$ with $\alpha-\beta=c$ or $\alpha-\beta=ci$.
\end{trybox}

\hsub{Implication $\Rightarrow$ (``if \dots\ then'')}
\begin{key}[{Definition 1.1.27 \textperiodcentered\ rules for $\Rightarrow$}]
($\Rightarrow$I) if $q$ can be derived from the assumption that $p$ is true, then $p\Rightarrow q$ is true;\quad ($\Rightarrow$E, \emph{modus ponens}) if $p\Rightarrow q$ and $p$ are true, then $q$ is true.
\end{key}
\begin{strategy}[{Strategies 1.1.28 and 1.1.34}]
\textbf{Proving $p\Rightarrow q$:} assume $p$, derive $q$. First words: ``Suppose $p$.'' \textbf{Assuming $p\Rightarrow q$:} if you also have $p$, conclude $q$. Used to reduce a goal: if $p\Rightarrow q$ is known and $p$ is easy, prove $p$ instead of $q$ (Example 1.1.36).
\end{strategy}
\begin{result}[{Proposition 1.1.30}]
Let $x,y\in\R$. If $x$ and $x+y$ are rational, then $y$ is rational. \emph{Proof:} suppose $x=\frac ab$, $x+y=\frac cd$; then $y=(x+y)-x=\frac{bc-ad}{bd}\in\Q$.
\end{result}

\begin{bookex}[essential]{1.1.32}
Let $r$ and $s$ be propositions. Prove that $r\Rightarrow(s\Rightarrow s)$ is true, using the strategies for implications.
\solution
By Strategy 1.1.28 it suffices to assume $r$ and derive $s\Rightarrow s$. So assume $r$. To prove $s\Rightarrow s$, again by Strategy 1.1.28 it suffices to assume $s$ and derive $s$. So assume $s$. Then $s$ is true, since it is assumed. Hence $s\Rightarrow s$ is true, and therefore $r\Rightarrow(s\Rightarrow s)$ is true.\qed
\end{bookex}

\begin{bookex}[essential]{1.1.33}
Let $a,b,c$ be propositions. Prove that $(a\vee b)\Rightarrow\big(c\Rightarrow((c\wedge a)\vee b)\big)$ is true.
\solution
Assume $a\vee b$ (Strategy 1.1.28). We must derive $c\Rightarrow((c\wedge a)\vee b)$, so assume $c$ as well. Now split into cases on $a\vee b$ (Strategy 1.1.19).\\
\emph{Case 1:} $a$ is true. Since $c$ and $a$ are both true, $c\wedge a$ is true (Strategy 1.1.7), and so $(c\wedge a)\vee b$ is true (Strategy 1.1.16).\\
\emph{Case 2:} $b$ is true. Then $(c\wedge a)\vee b$ is true (Strategy 1.1.16).\\
In both cases $(c\wedge a)\vee b$ holds. Hence $c\Rightarrow((c\wedge a)\vee b)$, and hence the whole implication, is true.\qed
\end{bookex}

\begin{bookex}[recommended]{1.1.37}
Let $p,q,r$ be propositions. Prove that $(p\Rightarrow(q\Rightarrow r))\Rightarrow(q\Rightarrow(p\Rightarrow r))$ is true.
\solution
Assume $p\Rightarrow(q\Rightarrow r)$. To prove $q\Rightarrow(p\Rightarrow r)$, assume $q$; to prove $p\Rightarrow r$, assume $p$. By modus ponens (Strategy 1.1.34) applied to $p\Rightarrow(q\Rightarrow r)$ and $p$, we get $q\Rightarrow r$; by modus ponens again with $q$, we get $r$. Hence $p\Rightarrow r$, hence $q\Rightarrow(p\Rightarrow r)$, hence the whole formula.\qed
\end{bookex}

\begin{bookex}[essential]{1.1.38}
Let $p,q,r$ be propositions. Prove that $(p\Rightarrow p\wedge q)\Rightarrow((p\vee q)\Rightarrow q)$ is true.
\solution
Assume $p\Rightarrow(p\wedge q)$. To prove $(p\vee q)\Rightarrow q$, assume $p\vee q$ and split into cases.\\
\emph{Case 1:} $p$ is true. By modus ponens, $p\wedge q$ is true, so $q$ is true by Strategy 1.1.10 ($\wedge$E$_2$).\\
\emph{Case 2:} $q$ is true. Then $q$ is true.\\
In both cases $q$ holds, so $(p\vee q)\Rightarrow q$ is true, and hence the whole formula is true.\qed
\end{bookex}

\begin{trybox}[{1.1.31 \fO}]
\textbf{1.1.31} Let $p(x)$ be a polynomial over $\C$. Prove that if $\alpha$ is a root of $p(x)$ and $a\in\C$, then $\alpha$ is a root of $(x-a)p(x)$.
\end{trybox}

\hsub{Converse and biconditional $\Leftrightarrow$}
\begin{key}[{Definitions 1.1.39 and 1.1.40}]
The \term{converse} of $p\Rightarrow q$ is $q\Rightarrow p$. Ways of saying $p\Rightarrow q$: ``if $p$ then $q$'', ``$p$ only if $q$'', ``$p$ is sufficient for $q$''. Ways of saying $q\Rightarrow p$: ``$p$ if $q$'', ``$p$ is necessary for $q$''. \\
$p\Leftrightarrow q$ means $(p\Rightarrow q)\wedge(q\Rightarrow p)$, read ``$p$ if and only if $q$''. So: \textbf{to prove $p\Leftrightarrow q$, prove both implications}, labelled ($\Rightarrow$) and ($\Leftarrow$); \textbf{to use $p\Leftrightarrow q$}, use either implication.
\end{key}
\begin{warn}[{Examples 1.1.44 and 1.1.45 \textperiodcentered\ solving an equation is a biconditional}]
Rearranging an equation only proves ``if $x$ is a solution then $x=\dots$'' (it narrows the candidates). You must also check the converse by substituting back: for $x+\sqrt x=0$ the rearrangement gives $x=0$ or $x=1$, but $1$ is not a solution.
\end{warn}
\begin{key}[{Example 1.1.46 \textperiodcentered\ divisibility tests}]
To prove ``$n$ is divisible by $8$ iff its last three digits form a number divisible by $8$'': write $n=n'+n''$ with $n'$ the last three digits and $n''=n-n'=1000\cdot(\dots)$, note $8$ divides $n''$, and use Exercise 0.40 in both directions ($n'=1\cdot n+(-1)\cdot n''$ and $n=1\cdot n'+1\cdot n''$).
\end{key}

\begin{bookex}[recommended]{1.1.43}
Let $p$ and $q$ be propositions. Prove that $(p\Leftrightarrow q)\Leftrightarrow(q\Leftrightarrow p)$ is true.
\solution
By Definition 1.1.40 we prove both implications.\\
($\Rightarrow$) Assume $p\Leftrightarrow q$, i.e.\ $(p\Rightarrow q)\wedge(q\Rightarrow p)$. By Strategy 1.1.10 we have $q\Rightarrow p$ and $p\Rightarrow q$, so by Strategy 1.1.7, $(q\Rightarrow p)\wedge(p\Rightarrow q)$ is true, which is $q\Leftrightarrow p$.\\
($\Leftarrow$) Assume $q\Leftrightarrow p$, i.e.\ $(q\Rightarrow p)\wedge(p\Rightarrow q)$. The same two steps give $(p\Rightarrow q)\wedge(q\Rightarrow p)$, i.e.\ $p\Leftrightarrow q$.\qed
\end{bookex}

\begin{bookex}[essential]{1.1.47}
Prove that a natural number $n$ is divisible by $3$ if and only if the sum of its base-10 digits is divisible by $3$.
\solution
Let $d_r\dots d_1d_0$ be the base-10 expansion of $n$, so $n=\sum_{i=0}^rd_i10^i$, and let $s=\sum_{i=0}^rd_i$ be the digit sum. Put $n''=n-s=\sum_{i=0}^rd_i(10^i-1)$.\\
\emph{Claim: $3$ divides $n''$.} For each $i$, $10^i-1=(10-1)(10^{i-1}+10^{i-2}+\dots+1)=9m_i$ with $m_i\in\N$ (for $i=0$, $10^0-1=0=9\cdot0$), so $10^i-1=3\cdot(3m_i)$. Hence $n''=3\sum_id_i\cdot3m_i$ is divisible by $3$.\\
($\Rightarrow$) Suppose $3$ divides $n$. Since $3$ divides $n''$, Exercise 0.40 gives that $3$ divides $1\cdot n+(-1)\cdot n''=n-n''=s$.\\
($\Leftarrow$) Suppose $3$ divides $s$. Since $3$ divides $n''$, Exercise 0.40 gives that $3$ divides $1\cdot s+1\cdot n''=n$.\qed
\end{bookex}

\hsub{Negation $\neg$ (``not'') and contradiction $\bot$}
\begin{key}[{Definitions 1.1.48 and 1.1.50}]
A \term{contradiction} is a proposition known or assumed to be false; $\bot$ denotes an arbitrary one (e.g.\ $0=1$, $\sqrt2\in\Q$). Rules for $\neg$: ($\neg$I) if a contradiction can be derived from the assumption that $p$ is true, then $\neg p$ is true; ($\neg$E) if $\neg p$ and $p$ are both true, a contradiction may be derived.
\end{key}
\begin{strategy}[{Strategies 1.1.51 and 1.1.59 \textperiodcentered\ proof by contradiction}]
\textbf{Proving $\neg p$:} assume $p$ and derive a contradiction. First words: ``Suppose $p$ (towards a contradiction).'' \textbf{Assuming $\neg p$:} any derivation of $p$ is a contradiction; this is how the contradiction is usually reached.\\
Statements that are negations in disguise: ``$x$ is irrational'' ($=\neg(x\in\Q)$), ``$b$ does not divide $a$'', ``there is no \dots'', ``$X$ is empty''.
\end{strategy}
\begin{result}[{Propositions 1.1.53 and 1.1.58}]
\textbf{1.1.53} If $r\in\Q$ and $a$ is irrational, then $r+a$ is irrational. \emph{Proof:} suppose $r+a\in\Q$; then $a=(r+a)-r\in\Q$ by Proposition 1.1.30, contradicting the assumption. \\
\textbf{1.1.58} An integer $a$ is odd iff $a=2b+1$ for some $b\in\Z$. \emph{Proof:} ($\Rightarrow$) by the division theorem $a=2b$ or $a=2b+1$; the first contradicts ``odd''. ($\Leftarrow$) $a=2b+1$ leaves remainder $1$, and even numbers leave remainder $0$, so $a$ is not even.
\end{result}

\begin{bookex}[recommended]{1.1.54}
Let $p$ and $q$ be propositions. Prove that $(p\Rightarrow\neg q)\Rightarrow(q\Rightarrow\neg p)$ is true.
\solution
Assume $p\Rightarrow\neg q$. To prove $q\Rightarrow\neg p$, assume $q$. To prove $\neg p$, we use Strategy 1.1.51: assume $p$ and derive a contradiction. By modus ponens from $p\Rightarrow\neg q$ and $p$ we obtain $\neg q$. But $q$ is also assumed, so by ($\neg$E) we have a contradiction. Hence $\neg p$ is true, so $q\Rightarrow\neg p$ is true, so the whole formula is true.\qed
\end{bookex}

\begin{bookex}[recommended]{1.1.55}
Let $p$ and $q$ be propositions. Prove that $\neg p\Rightarrow(\neg q\Rightarrow\neg(p\vee q))$ is true.
\solution
Assume $\neg p$ and assume $\neg q$. To prove $\neg(p\vee q)$, assume $p\vee q$ and derive a contradiction, by cases. \emph{Case 1:} $p$. This contradicts $\neg p$. \emph{Case 2:} $q$. This contradicts $\neg q$. In both cases we have a contradiction, so $\neg(p\vee q)$ is true. Hence $\neg q\Rightarrow\neg(p\vee q)$, and hence the whole formula, is true.\qed
\end{bookex}

\begin{bookex}[essential]{1.1.56}
Let $x\in\R$. Prove by contradiction that if $x$ is irrational then $-x$ and $\frac1x$ are irrational.
\solution
Suppose $x$ is irrational. (Note $x\ne0$, since $0=\frac01$ is rational.)\\
\emph{$-x$ is irrational.} Suppose, towards a contradiction, that $-x$ is rational, say $-x=\frac ab$ with $a,b\in\Z$, $b\ne0$. Then $x=\frac{-a}{b}$ with $-a\in\Z$, so $x$ is rational, contradicting the assumption. Hence $-x$ is irrational.\\
\emph{$\frac1x$ is irrational.} Suppose $\frac1x=\frac ab$ with $a,b\in\Z$, $b\ne0$. Since $\frac1x\ne0$ we have $a\ne0$, and then $x=\frac ba$ with $b,a\in\Z$ and $a\ne0$, so $x$ is rational: contradiction. Hence $\frac1x$ is irrational.\qed
\end{bookex}

\begin{bookex}[essential]{1.1.57}
Prove by contradiction that there is no least positive real number: there is no positive real $a$ such that $a\le b$ for all positive real $b$.
\solution
Suppose, towards a contradiction, that $a$ is a positive real number such that $a\le b$ for every positive real number $b$. Then $b=\frac a2$ is a positive real number, so $a\le\frac a2$. Multiplying by $2$ gives $2a\le a$, hence $a\le0$, contradicting $a>0$. So no such $a$ exists.\qed
\end{bookex}

\hsub{Two logical axioms}
\begin{result}[{Axioms 1.1.60 and 1.1.67}]
\textbf{Law of excluded middle (LEM):} for every propositional formula $p$, the formula $p\vee(\neg p)$ is true. \textbf{Principle of explosion:} from a contradiction, any consequence may be derived.
\end{result}
\begin{strategy}[{Strategy 1.1.61 \textperiodcentered\ using LEM}]
To prove $q$, you may split into the cases ``$p$ true'' and ``$p$ false'' for any proposition $p$ you like, and prove $q$ in each case. (Proposition 1.1.63: if $ab$ is even then $a$ or $b$ is even; split on ``$a$ even''.)
\end{strategy}

\begin{bookex}[recommended]{1.1.65}
Let $p$ and $q$ be propositions. Prove that $(\neg q\Rightarrow\neg p)\Rightarrow(p\Rightarrow q)$ is true.
\solution
Assume $\neg q\Rightarrow\neg p$, and assume $p$; we must derive $q$. By the law of excluded middle, $q$ or $\neg q$.\\
\emph{Case 1:} $q$ is true. Done.\\
\emph{Case 2:} $\neg q$ is true. By modus ponens, $\neg p$ is true. Together with $p$ this is a contradiction ($\neg$E), from which $q$ follows by the principle of explosion.\\
In both cases $q$ holds. Hence $p\Rightarrow q$, and hence the whole formula, is true.\qed
\end{bookex}

\begin{bookex}[recommended]{1.1.66}
Let $a$ and $b$ be irrational. By considering $\sqrt2^{\sqrt2}$, prove that it is possible that $a^b$ be rational.
\solution
By the law of excluded middle, $\sqrt2^{\sqrt2}$ is rational or it is not.\\
\emph{Case 1:} $\sqrt2^{\sqrt2}$ is rational. Take $a=b=\sqrt2$, both irrational (Assumption 0.49); then $a^b=\sqrt2^{\sqrt2}$ is rational.\\
\emph{Case 2:} $\sqrt2^{\sqrt2}$ is irrational. Take $a=\sqrt2^{\sqrt2}$ and $b=\sqrt2$, both irrational. Then $a^b=\big(\sqrt2^{\sqrt2}\big)^{\sqrt2}=\sqrt2^{\sqrt2\cdot\sqrt2}=\sqrt2^{\,2}=2$, which is rational.\\
In both cases we found irrational $a,b$ with $a^b$ rational.\qed\ (We never learn which case holds: a \emph{non-constructive} proof, which is why LEM is called non-constructive.)
\end{bookex}

\begin{trybox}[{1.1.64 \fO}]
\textbf{1.1.64} In the proof of Proposition 1.1.63, where is LEM used, where is contradiction used, and what is the contradiction? Prove the proposition again twice, once without LEM and once without contradiction.
\end{trybox}

\begin{warn}
\begin{itemize}
\item Keep the direction straight: to prove $p\Rightarrow q$ assume $p$; proving $q\Rightarrow p$ instead (the converse) proves a different proposition (1.E.12).
\item When you use the law of excluded middle or proof by cases, the cases must be exhaustive and \emph{every} case must reach the goal.
\item In a proof by contradiction, say what the contradiction is: ``this contradicts the assumption that $x$ is irrational''.
\end{itemize}
\end{warn}

\hsection{Variables and quantifiers}

\begin{key}[{Definitions 1.2.1, 1.2.2, 1.2.4}]
In a statement, a variable $x$ (understood to range over a set $X$) is \term{free} if it makes sense to substitute elements of $X$ for it, and \term{bound} otherwise (``there exist integers $a,b$ with $x=a^2+b^2$'': $x$ free, $a,b$ bound). A \term{predicate} $p(x_1,\dots,x_n)$ is a symbol with a list of free variables, each with a \term{domain of discourse}; it becomes a proposition when values are substituted. A \term{logical formula} is built from predicates with logical operators and the \term{quantifiers} $\forall$ (``for all $x\in X$''), $\exists$ (``there exists $x\in X$ such that''), $\exists!$ (``there is exactly one''). Quantifiers go \emph{in front}: ``$n=2k$ for some $k\in\Z$'' is written $\exists k\in\Z,\ n=2k$.
\end{key}

\begin{key}[{Translating (Examples 1.2.5, 1.2.7)}]
English $\to$ symbols: introduce a variable for each noun phrase, unfold definitions (``$n$ is even'' is $\exists k\in\Z,\ n=2k$), and restrict with an implication inside $\forall$ (``every positive real has a square root'': $\forall a\in\R,\ (a>0\Rightarrow\exists b\in\R,\ b^2=a)$). Symbols $\to$ English: read symbol by symbol, then remove bound variables by naming the concept they encode (``$\exists b\in\R,\ a=b^2$'' is ``$a$ has a real square root'').
\end{key}

\begin{bookex}[essential]{1.2.6}
Find logical formulae representing: (a) there is an integer that is divisible by every integer; (b) there is no greatest odd integer; (c) between any two distinct rational numbers is a third distinct rational number; (d) if an integer has a rational square root, then that root is an integer.
\solution
(a) $\exists n\in\Z,\ \forall m\in\Z,\ \exists q\in\Z,\ n=qm$. \reason{$m$ divides $n$ unfolded}\\
(b) $\neg\exists n\in\Z,\ \big([\exists k\in\Z,\ n=2k+1]\wedge\forall m\in\Z,\ ([\exists\ell\in\Z,\ m=2\ell+1]\Rightarrow m\le n)\big)$; equivalently $\forall n\in\Z,\ \big(n\text{ odd}\Rightarrow\exists m\in\Z,\ (m\text{ odd}\wedge m>n)\big)$.\\
(c) $\forall a\in\Q,\ \forall b\in\Q,\ \big(a<b\Rightarrow\exists c\in\Q,\ (a<c\wedge c<b)\big)$.\\
(d) $\forall n\in\Z,\ \forall x\in\Q,\ (x^2=n\Rightarrow x\in\Z)$.\qed
\end{bookex}

\begin{trybox}[{1.2.8 \fO}]
\textbf{1.2.8} Plain English, as few variables as possible: (a) $\exists q\in\Z,\ a=qb$; (b) $\exists a\in\Z,\ \exists b\in\Z,\ (b\ne0\wedge bx=a)$; (c) $\forall d\in\N,\ [(\exists q\in\Z,\ n=qd)\Rightarrow(d=1\vee d=n)]$; (d) $\forall a\in\R,\ [a>0\Rightarrow\exists b\in\R,\ (b>0\wedge a<b)]$.
\end{trybox}

\hsub{Universal quantifier $\forall$}
\begin{key}[{Definition 1.2.9 \textperiodcentered\ rules for $\forall$}]
($\forall$I) if $p(x)$ can be derived from the assumption that $x$ is an \emph{arbitrary} element of $X$, then $\forall x\in X,\ p(x)$;\quad ($\forall$E) if $\forall x\in X,\ p(x)$ is true and $a\in X$, then $p(a)$ is true.
\end{key}
\begin{strategy}[{Strategies 1.2.10 and 1.2.19}]
\textbf{Proving $\forall x\in X,\ p(x)$:} ``Let $x\in X$'' (or ``fix'', ``take''), then prove $p(x)$ assuming nothing about $x$ except $x\in X$. \textbf{Assuming $\forall x\in X,\ p(x)$:} you may use $p(a)$ for any particular $a\in X$ you need.
\end{strategy}
\begin{warn}[{Common error (page 54)}]
``Let $n$ be an arbitrary integer, say $n=17$'' proves $\exists n$, not $\forall n$. The \emph{reader} chooses $n$, not you. A correct proof of $\forall n\in\Z,\ n^2\ge0$: let $n\in\Z$; either $n\ge0$ or $n<0$; in each case $n^2\ge0$.
\end{warn}
\begin{result}[{Propositions 1.2.11 and 1.2.12}]
The square of every odd integer is odd ($n=2k+1\Rightarrow n^2=2(2k^2+2k)+1$). The last decimal digit of a square is $0,1,4,5,6$ or $9$ (write $n=10m+d_0$; then $n^2=10m(10m+2d_0)+d_0^2$ ends in the last digit of $d_0^2$).
\end{result}

\begin{bookex}[essential]{1.2.15}
Prove that for every integer $n$, $n$ is even if and only if $n^2$ is even.
\solution
Let $n\in\Z$. We prove both directions.\\
($\Rightarrow$) Suppose $n$ is even, so $n=2k$ for some $k\in\Z$. Then $n^2=4k^2=2(2k^2)$ with $2k^2\in\Z$, so $n^2$ is even.\\
($\Leftarrow$) Suppose $n^2$ is even. By the division theorem, $n=2k$ or $n=2k+1$ for some $k\in\Z$. If $n=2k+1$, then $n^2$ is odd by Proposition 1.2.11, contradicting that $n^2$ is even. So $n=2k$, i.e.\ $n$ is even.\qed
\end{bookex}

\begin{bookex}[essential]{1.2.18}
Prove that for all real numbers $x$ and $y$, if $x$ is irrational, then $x+y$ and $x-y$ are not both rational.
\solution
Let $x,y\in\R$ and suppose $x$ is irrational. Suppose, towards a contradiction, that $x+y$ and $x-y$ are both rational. Then by Theorem 0.16 their sum $(x+y)+(x-y)=2x$ is rational, and hence $x=\frac{2x}2$ is rational (Theorem 0.16 again, dividing by $2\ne0$). This contradicts the assumption that $x$ is irrational. Hence $x+y$ and $x-y$ are not both rational.\qed
\end{bookex}

\begin{bookex}[recommended]{1.2.21 and 1.2.22}
\textbf{1.2.21} Suppose $\forall x\in\R,\ P(x)$. Show that we can conclude $P(21)$.\quad \textbf{1.2.22} Let $P(x),Q(x)$ be predicates with $x\in\R$. Prove that $(\forall x\in\R,\ P(x)\Rightarrow Q(x))\Rightarrow(\forall x\in\R,\ P(x))\Rightarrow(\forall x\in\R,\ Q(x))$.
\solution
\textbf{1.2.21} Since $21\in\R$ and $\forall x\in\R,\ P(x)$ holds, the rule ($\forall$E) (Strategy 1.2.19) gives $P(21)$.\qed\\
\textbf{1.2.22} Assume $\forall x\in\R,\ (P(x)\Rightarrow Q(x))$ and assume $\forall x\in\R,\ P(x)$. To prove $\forall x\in\R,\ Q(x)$, let $x\in\R$ be arbitrary. By ($\forall$E), $P(x)\Rightarrow Q(x)$ and $P(x)$ are both true, so $Q(x)$ is true by modus ponens. Since $x$ was arbitrary, $\forall x\in\R,\ Q(x)$.\qed
\end{bookex}

\begin{trybox}[{1.2.14 (bonus), 1.2.16, 1.2.17 \fO}]
\textbf{1.2.14} Prove $\forall x\in\R,\ x=x$.\quad \textbf{1.2.16} Prove that every integer is rational.\quad \textbf{1.2.17} Prove that every linear polynomial over $\Q$ has a rational root.
\end{trybox}

\hsub{Existential quantifier $\exists$}
\begin{key}[{Definition 1.2.23 \textperiodcentered\ rules for $\exists$}]
($\exists$I) if $a\in X$ and $p(a)$ is true, then $\exists x\in X,\ p(x)$;\quad ($\exists$E) if $\exists x\in X,\ p(x)$ is true, and $q$ can be derived from the assumption that $p(a)$ holds for some fixed $a\in X$, then $q$ is true.
\end{key}
\begin{strategy}[{Strategies 1.2.24 and 1.2.30}]
\textbf{Proving $\exists x\in X,\ p(x)$:} define a specific $a$ (``Define $n=9$'', ``Take $z=\frac{x+y}2$''), check $a\in X$, and prove $p(a)$. \textbf{Assuming $\exists x\in X,\ p(x)$:} ``Take $a\in X$ such that $p(a)$'' with a letter not yet in use. Divisibility assumptions are existential: from ``$3$ divides $n^3$'' take $q$ with $n^3=3q$ (Proposition 1.2.31).
\end{strategy}

\begin{bookex}[recommended]{1.2.28}
Prove that there is an integer which is divisible by zero.
\solution
Define $n=0$. Then $n\in\Z$, and $0=0\cdot0$ with $0\in\Z$, so $0$ divides $n$ by Definition 0.36. Hence $\exists n\in\Z$ such that $0$ divides $n$.\qed
\end{bookex}

\begin{bookex}{1.2.29}
Prove that for all $x,y\in\Q$, if $x<y$ then there is some $z\in\Q$ with $x<z<y$.
\solution
Let $x,y\in\Q$ and suppose $x<y$. Define $z=\frac{x+y}2$. Then $z\in\Q$ by Theorem 0.16 ($x+y\in\Q$ and $2\ne0$). Moreover $x<z$ because $2x<x+y$ (add $x$ to $x<y$), and $z<y$ because $x+y<2y$. So $x<z<y$, and $z$ is the required element.\qed
\end{bookex}

\begin{trybox}[{1.2.27 \fO}]
\textbf{1.2.27} Prove that there is a real number that is irrational but whose square is rational.
\end{trybox}

\hsub{Unique existence $\exists!$}
\begin{key}[{Definition 1.2.32 and Proposition 1.2.38}]
$\exists!x\in X,\ p(x)$ is shorthand for
\[
\underbrace{(\exists x\in X,\ p(x))}_{\text{existence}}\ \wedge\ \underbrace{(\forall a\in X,\ \forall b\in X,\ [(p(a)\wedge p(b))\Rightarrow a=b])}_{\text{uniqueness}} ,
\]
equivalently (1.2.38) $\exists x\in X,\ \big(p(x)\wedge\forall y\in X,\ (p(y)\Rightarrow y=x)\big)$.
\end{key}
\begin{strategy}[{Strategy 1.2.35}]
\textbf{Existence:} exhibit $a$ with $p(a)$. \textbf{Uniqueness:} ``Let $a,b\in X$ with $p(a)$ and $p(b)$'' \dots\ ``so $a=b$''. (Or: with the $a$ from existence, show $\forall y,\ p(y)\Rightarrow y=a$.) Example 1.2.36: for $a>0$ the positive square root is unique because $y^2=z^2$ gives $(y-z)(y+z)=0$ and $y+z>0$.
\end{strategy}

\begin{bookex}[recommended]{1.2.37 (a), (b)}
Write as a formula with $\exists!$ and prove: (a) for each real $a$, the equation $x^2+2ax+a^2=0$ has exactly one real solution; (b) there is a unique real $a$ for which $x^2+a^2=0$ has a real solution.
\solution
(a) $\forall a\in\R,\ \exists!x\in\R,\ x^2+2ax+a^2=0$. Let $a\in\R$. \emph{Existence:} $x=-a$ gives $a^2-2a^2+a^2=0$. \emph{Uniqueness:} let $x,y\in\R$ both be solutions. Since $x^2+2ax+a^2=(x+a)^2$, we get $(x+a)^2=0$, so $x+a=0$, so $x=-a$; likewise $y=-a$; hence $x=y$.\qed\\
(b) $\exists!a\in\R,\ \exists x\in\R,\ x^2+a^2=0$. \emph{Existence:} $a=0$ works, with $x=0$. \emph{Uniqueness:} let $a,b\in\R$ both have the property; take $x\in\R$ with $x^2+a^2=0$. Since $x^2\ge0$ and $a^2\ge0$, the sum is $0$ only if $a^2=0$, so $a=0$. The same argument gives $b=0$, so $a=b$.\qed
\end{bookex}

\begin{trybox}[{1.2.37 (c), 1.2.34 \fO}]
\textbf{1.2.37 (c)} There is a unique natural number with exactly one positive divisor.\quad \textbf{1.2.34} (Discussion) Why does Definition 1.2.32 capture ``exactly one''? Other possible definitions?
\end{trybox}

\hsub{Quantifier alternation}
\begin{result}[{Theorem 1.2.40}]
$\exists y\in Y,\ \forall x\in X,\ p(x,y)\ \Rightarrow\ \forall x\in X,\ \exists y\in Y,\ p(x,y)$. \emph{Proof:} fix $a\in X$; choose $b$ with $\forall x,\ p(x,b)$; then $p(a,b)$. \\
The converse fails in general: in $\forall x\,\exists y$ the $y$ may depend on $x$ (``every door has a key''); in $\exists y\,\forall x$ one $y$ works for all $x$ (``a master key'').
\end{result}

\begin{bookex}[essential]{1.2.39}
Let $p(x,y)$ be ``$x+y$ is even''. Prove that $\forall x\in\Z,\ \exists y\in\Z,\ p(x,y)$ is true, and that $\exists y\in\Z,\ \forall x\in\Z,\ p(x,y)$ is false.
\solution
\emph{First statement.} Let $x\in\Z$. Define $y=x$. Then $y\in\Z$ and $x+y=2x$ is even (it is $2\cdot x$ with $x\in\Z$). So $\exists y\in\Z,\ p(x,y)$, and as $x$ was arbitrary the statement holds.\\
\emph{Second statement.} Suppose, towards a contradiction, that there is $y\in\Z$ with $x+y$ even for every $x\in\Z$. Take $x=1-y\in\Z$. Then $x+y=1$, so $1$ is even, i.e.\ $1=2k$ for some $k\in\Z$; but then $k=\frac12\notin\Z$, a contradiction. Hence the statement is false.\qed
\end{bookex}
